\thispagestyle{plain}
\begin{center}
{\fontsize{10}{12}\selectfont Manuscrit de compilation pour le transfert académique\par}
\vspace{6mm}
{\fontsize{17}{20.5}\selectfont\bfseries Arithmétique généalogique des nombres premiers et structure spectrale de la fonction zêta de Riemann\par}
\vspace{7pt}
{\fontsize{12}{14.5}\selectfont Transport arithmétique, moments spectraux et critère de positivité de Weil\par}
\vspace{8pt}
{\fontsize{11.5}{14}\selectfont Oumar Haidara Fall\space\textperiodcentered\space Rubén Ramos Balsa\par}
\vspace{4pt}
{\fontsize{8}{10}\selectfont Coautorat sans préséance de contribution\par}
\vspace{4pt}
{\fontsize{8}{10}\selectfont Intégration et compilation documentaires générées par ChatGPT - Astra.\par}
\end{center}
\vspace{3pt}
\phantomsection
\addcontentsline{toc}{section}{Résumé}
\begin{center}\bfseries Résumé\end{center}
\begingroup\fontsize{10.5}{12.8}\selectfont
On démontre un critère équivalent de positivité de la forme arithmétique de Weil, dans chaque fenêtre fixée, au moyen d'un complément de Schur de taille finie qui conserve l'effet de l'espace infini des détails. Son application établit $\mathscr W(f,f)\geq\frac12\|f\|_2^2$ sur le domaine conjoint des lecteurs supporté dans $(-1/18,1/18)$, ainsi que dans ses translations et raffinements intérieurs. On construit en outre un lecteur qui récupère la fonction test à partir de la queue gamma et des transports réversibles préservant la forme complète, y compris les contributions première, gamma et polaire. La borne se transmet à des collections entières de fonctions tests à queues exponentielles.

Ces résultats sont développés à partir d’un noyau formel mathématique
dénommé holographie modulaire triadique, ci-après HMT.
\emph{All Possible Paths} (APP) est la structure arithmétique de
chemins sur le support toroïdal $(\mathbb Z/9\mathbb Z)^2$,
dont les translations cardinales forment un graphe de Cayley additif ;
ses feuillets conservent somme et produit avec reste et quotient. TRIT est
la signature orientée $\{-1,0,+1\}$ qui distingue les régimes
locaux. Le \emph{Topological Phase Kernel} (TPK) compose
sélection, transport, retour et actualisation de la mémoire sur
l’état enrichi. Les réalisations arithmétique et spectrale
utilisent la structure discrète conjointe du continu produite
par ce développement.

Le produit natif sur les mots nonadiques possède une forme normale unique ; ses ouvertures multiplicatives sélectionnent les irréductibles que l'évaluation identifie aux nombres premiers. Les occupations de ces modes génèrent le produit d'Euler et les moments premier--puissance. La limite des traces de chaleur cycliques et la transformation de Mellin produisent la fonction complétée et son équation fonctionnelle. La reconstruction de l'intérieur d'APP à partir de la frontière relevée, les feuillets de continuation et les vacances précisent quelle information conserve le raffinement.

Le double cercle relie la coordonnée spectrale à la mémoire des retours. La compression nonadique détermine le bilan $8/81$ et les modes centraux le rapport $5/9$, dont les moments admettent une factorisation positive en toute dimension. Pour la forme arithmétique, le contrôle des détails et des termes croisés fournit la réduction de Schur et la coercivité indiquée. Le critère global de Weil est formulé au moyen de l'identification conjointe des moments arithmétiques à un Gram positif, y compris tous les moments mixtes et la couverture du domaine des fonctions tests. Cette identification globale demeure une hypothèse des résultats spectraux qui l'utilisent : les preuves de coercivité établies ici ne constituent pas une démonstration de l'hypothèse de Riemann.

\medskip
\noindent\textbf{Mots-clés :} forme de Weil ; coercivité ;
complément de Schur ; nombres premiers ; fonction zêta de Riemann ;
holographie modulaire triadique ; produit natif ; reconstruction de
frontière ; mémoire ; moments mixtes.
\par\endgroup
